\begin{onehalfspace}

I Capitoli~\ref{chap:qec} e~\ref{chap:surface} studiano la correzione d'errore su circuiti dedicati, mentre il Capitolo~\ref{chap:risultati} misura la degradazione di Shor sotto un modello di rumore fisico. Collegare direttamente i due insiemi di risultati richiederebbe una compilazione fault-tolerant completa: codifica dei registri, estrazione ripetuta della sindrome, decoder, operazioni logiche e distillazione delle risorse non-Clifford. Questa tesi non simula tale architettura. Il presente capitolo svolge invece un esperimento di \textbf{sensibilit\`a fenomenologica}: applica al circuito compilato un errore Pauli indipendente per gate, indicato con $p_L$, e misura come varia la probabilit\`a di fattorizzazione. Il pedice ``L'' identifica il ruolo concettuale del parametro, ma non autorizza a sostituirvi direttamente le metriche logiche aggregate dei capitoli precedenti.

% -----------------------------------------------
\section{L'architettura a quattro livelli}
\label{sec:quattro_livelli}

La parte sperimentale separa quattro oggetti: Shor ideale, Shor sotto rumore fisico, benchmark QEC isolati e proxy Pauli per gate compilato. La separazione è necessaria perché Stim/PyMatching caratterizza efficientemente circuiti stabilizer, mentre Shor contiene aritmetica e rotazioni non-Clifford; sostituire i suoi dodici qubit con blocchi codificati, ancille e cicli di sindrome richiederebbe una simulazione fault-tolerant completa. Il quarto livello risponde quindi a una domanda di sensibilità e non riceve automaticamente i tassi aggregati del terzo.

% -----------------------------------------------
\section{Definizione operativa del proxy}
\label{sec:modello_logico}

In M8, $p_L$ \`e definito come la \textbf{probabilit\`a totale di applicare un Pauli non identit\`a dopo ciascun gate compilato incluso nel modello}. Per un gate su $q$ qubit il canale distribuisce $p_L$ uniformemente sui $4^q-1$ Pauli non identit\`a. Poich\'e Qiskit Aer parametrizza il canale depolarizzante con $\lambda_q$, lo script usa
\begin{equation}
    \lambda_1=\frac{4}{3}p_L,
    \qquad
    \lambda_2=\frac{16}{15}p_L.
\end{equation}
Il rumore \`e applicato a \texttt{sx}/\texttt{x} e \texttt{cx}; le rotazioni \texttt{rz} sono considerate virtuali. Si tratta quindi di un proxy uniforme \emph{per gate nella base compilata}, non di una probabilit\`a per ciclo di correzione, per round di sindrome o per intero blocco logico.

Il circuito è quello canonico $N=15$, $a=7$, con otto qubit di conteggio e quattro di lavoro, compilato deterministicamente nella base \texttt{rz/sx/x/cx}. Ogni punto M8 aggrega 20 repliche da 4096 shot; il successo vale uno quando frazioni continue e massimo comun divisore restituiscono fattori non banali. Manifest e hash della netlist fissano l'identità del circuito.

Questa definizione rende riproducibile la domanda sperimentale:
\begin{quote}
\emph{come cambia il successo di questa specifica implementazione di Shor al crescere di un errore Pauli indipendente e uniforme per gate?}
\end{quote}
Non consente invece di rispondere, senza un ulteriore modello architetturale, a quale distanza di surface code corrisponda un dato punto della curva.

% -----------------------------------------------
\section{Il contesto delle stime fault-tolerant}
\label{sec:resource_estimates}

Le stime crittografiche includono routing, cicli di sindrome, operazioni logiche e fabbriche di stati magici; non si ricavano estrapolando $N=15$~\cite{gidney2021}. La curva seguente è quindi locale al circuito e al canale dichiarati, non una previsione di risorse hardware.

% -----------------------------------------------
\section{Risultati M8 e M13}
\label{sec:risultati_logico}

\subsection*{M8: successo per singola misura}

La griglia preregistrata copre $p_L=0$--$0.5$. Gli intervalli sono Wilson al 95\%; la monotonia è controllata con copertura simultanea di Bonferroni sui confronti adiacenti.

\begin{figure}[H]
\centering
\includegraphics[width=0.82\textwidth]{qec_shor_logico_curve.pdf}
\caption[Shor sotto un proxy fenomenologico per gate]{Successo per singola misura in funzione dell'errore Pauli per gate compilato. Le metriche aggregate dei benchmark QEC non sono proiettate sulla curva perché non hanno la stessa unità.}
\label{fig:qec_shor_logico}
\end{figure}

La curva è compatibile con una funzione non crescente, ma non ogni passo adiacente è significativo. Il successo scende da $0.7489$ a $p_L=0$ a $0.4518$ per $p_L=0.02$ e raggiunge, oltre $0.1$, un plateau vicino a $0.245$. Quest'ultimo dipende dagli esiti che il post-processing di $N=15$ accetta e non è universale.

\subsection*{Ripetizioni indipendenti}

Il successo per misura non è il successo dopo più tentativi indipendenti. Se ciascun tentativo ha probabilità $P_s$,
\begin{equation}
    P(k)=1-(1-P_s)^k.
    \label{eq:cumulativa}
\end{equation}
Due tentativi portano così la stima a $0.937$ nel caso ideale e a circa $0.811$ per $p_L=0.01$. Una soglia dell'80\% è pertanto cumulativa, non un requisito sul singolo run.

\subsection*{M13: TOP-4 con pareggi espliciti}

M13 separa successo per misura, TOP-1 e TOP-4 sullo stesso istogramma. TOP-4 resta $1.000$ fino a $p_L=0.05$, ma scende a $0.915$ per $p_L=0.1$ e a $0.670$ per $p_L=0.2$: recupera candidati su questa piccola istanza, senza certificare la conservazione dello stato. Protocollo, tabella completa e analisi dei pareggi sono nella Sezione~\ref{app:compilazione} dell'Appendice~\ref{app:fase2}.

% -----------------------------------------------
\subsection{Discussione}
\label{subsec:discussione_logico}

M8 mostra che ridurre l'errore Pauli per gate riporta il successo verso il limite ideale dell'istanza. Non stabilisce quale codice produca quel parametro: errore Steane, errore del memory experiment surface-code e guadagno del decoder sono grandezze aggregate diverse e non possono essere inserite direttamente sulla curva. Il passaggio richiederebbe gate logici, numero di cicli per operazione, decoder e canali residui espliciti, comprese le operazioni non-Clifford.

$N=15$ resta inoltre un'istanza didattica, nella quale TOP-K beneficia di pochi picchi verificabili classicamente. Il capitolo caratterizza quindi la sensibilità del circuito controllato; non dimostra la fattibilità fault-tolerant di Shor su scala crittografica.

\end{onehalfspace}
